\documentclass[12pt]{article}
\usepackage[T1]{fontenc}
\usepackage[utf8]{inputenc}
\usepackage{lmodern}
\usepackage{textcomp}
\usepackage{enumitem}
\usepackage{geometry}
\usepackage{graphicx}
\usepackage{amsmath, amssymb}
\geometry{margin=1in}
\begin{document}
\begin{center}
History - K-12 Level Assessment 9 \
Total Marks: 40 \
Answer all questions.
\end{center}

\section*{Multiple Choice (10 marks)}
\begin{enumerate}[label=\arabic*.]
\item Which of the following characteristics of modern humans is the most ancient, pre-dating the Late Stone Age, the Upper Paleolithic and even anatomically modern humans?
\begin{enumerate}[label=(\alph*)]
    \item elaborate burials
    \item larger sites with larger populations
    \item improved stone tool technologies
    \item symbolic expression through the production of art
\end{enumerate}
\item Which of the following is true about the first apes?
\begin{enumerate}[label=(\alph*)]
    \item There were a far greater number of genera than today.
    \item They flourished in a broader area of the world than today's apes.
    \item Some ancient species were larger than today's apes.
    \item All of the above.
\end{enumerate}
\item Who was part of the elite class in the Olmec culture?
\begin{enumerate}[label=(\alph*)]
    \item those who lived in the particularly rich farmlands
    \item those who lived near resource-privileged areas
    \item those who lived near the river where trading took place
    \item those who lived near the volcanic rocks used in tool making
\end{enumerate}
\item Charles Darwin discovered that plants and animals had evolved through a process known as \_\_\_\_\_\_\_\_\_\_\_\_. Early archaeologists discovered that \_\_\_\_\_\_\_\_\_\_\_\_\_\_\_\_\_\_\_\_\_\_.
\begin{enumerate}[label=(\alph*)]
    \item scientific creationism; human culture evolved differently than plants and animals
    \item unilineal evolution; civilizations evolved more quickly than plants
    \item natural selection; human culture did not evolve or change much over time
    \item natural selection; human culture had evolved over an enormous period of time
\end{enumerate}
\item What distinguishes the trade patterns of the Mesolithic from those of the Upper Paleolithic?
\begin{enumerate}[label=(\alph*)]
    \item in the Mesolithic, goods traveled farther and more goods were exchanged
    \item in the Mesolithic, goods traveled farther but fewer goods were exchanged
    \item in the Mesolithic, goods traveled less far but more goods were exchanged
    \item in the Mesolithic, goods traveled less far and fewer goods were exchanged
\end{enumerate}
\end{enumerate}

\section*{True / False (10 marks)}
\textit{Answer True or False}
\begin{enumerate}[label=\arabic*.]
\setcounter{enumi}{5}
\item True or False: The ability to adapt to different environments was a key factor in the survival and success of anatomically modern humans.
\item True or False: The Aurignacian tool technology is associated with anatomically modern Homo sapiens.
\item True or False: Ancient South American civilizations created extensive written languages that documented their history.
\item True or False: During the Holocene epoch, Australia independently domesticated wheat and cattle, like some other parts of the world.
\item True or False: An activity area contains primary refuse, while a recycling area contains secondary refuse.
\end{enumerate}

\section*{Long Form Response (20 marks)}
\textit{Answer each question with a clear and complete explanation.}
\begin{enumerate}[label=\arabic*.]
\setcounter{enumi}{10}
\item Discuss the significance of parietal art in early human societies, using cave paintings as a primary example. What do they reveal about the culture and beliefs of their creators?
\item Discuss the significance of trace element analysis techniques in historical research, highlighting their impact on our understanding of past cultures and artifacts.
\end{enumerate}

\vfill
\noindent\textit{End of Paper}
\end{document}